\documentclass[11pt,a4paper]{article}

% ---------------------------------------------------------------
% OpenGAP — GitAgentProtocol: A Git-Native Protocol for the AI Agent Lifecycle
% arXiv preprint, cs.SE / cs.AI
% ---------------------------------------------------------------

\usepackage[margin=1in]{geometry}
\usepackage{amsmath,amssymb}
\usepackage{graphicx}
\usepackage{booktabs}
\usepackage{array}
\usepackage{longtable}
\usepackage{xcolor}
\usepackage{listings}
\usepackage[hidelinks]{hyperref}
\usepackage{authblk}
\usepackage{microtype}
\usepackage{caption}

\captionsetup{font=small,labelfont=bf}

\definecolor{codegray}{gray}{0.95}
\definecolor{codekey}{rgb}{0.13,0.29,0.53}
\definecolor{codestr}{rgb}{0.58,0.00,0.13}
\definecolor{codecom}{rgb}{0.30,0.45,0.30}

\lstdefinelanguage{yamlex}{
  keywords={true,false,null,yes,no},
  keywordstyle=\color{codekey}\bfseries,
  sensitive=false,
  comment=[l]{\#},
  commentstyle=\color{codecom}\itshape,
  morestring=[b]',
  morestring=[b]",
  stringstyle=\color{codestr}
}

\lstset{
  basicstyle=\ttfamily\footnotesize,
  backgroundcolor=\color{codegray},
  frame=single,
  rulecolor=\color{black!20},
  breaklines=true,
  columns=fullflexible,
  keepspaces=true,
  showstringspaces=false,
  aboveskip=6pt,
  belowskip=6pt,
  xleftmargin=6pt,
  xrightmargin=6pt
}

\title{\textbf{\Huge OpenGAP} \\[4pt] \large\textbf{GitAgentProtocol} \\[4pt] \normalsize\textit{A Git-Native Protocol for the AI Agent Lifecycle}}

\author[1]{Shreyas Kapale}
\affil[1]{OpenGAP Working Group \\ \texttt{shreyas@lyzr.ai} \\ \href{https://gitagent.sh}{gitagent.sh} \\ \href{https://github.com/open-gitagent/opengap}{github.com/open-gitagent/opengap}}

\date{\today}

\begin{document}

\maketitle

\begin{abstract}
\noindent
We present the \textbf{GitAgentProtocol} (OpenGAP), an open standard that treats a
git repository as the canonical unit of an AI agent. Where existing proposals such as
Anthropic's Model Context Protocol (MCP) and Google's Agent-to-Agent (A2A) standardize
the \emph{interfaces} between agents and the tools or peers they talk to, OpenGAP
standardizes the \emph{definition, lifecycle, and governance} of the agent itself:
identity (\texttt{SOUL.md}), constraints (\texttt{RULES.md}), segregation-of-duties
(\texttt{DUTIES.md}), declarative capabilities (\texttt{skills/}, \texttt{tools/},
\texttt{knowledge/}, \texttt{memory/}, \texttt{hooks/}), and regulatory-compliance
artifacts (\texttt{compliance/})\,---\,all as files under version control.

From this single source of truth, the reference CLI (\texttt{opengap}) deterministically
exports the same agent into fifteen different execution environments, including Claude
Code, OpenAI Agents SDK, CrewAI, Cursor, Gemini CLI, Codex, OpenCode, Kiro, Lyzr,
OpenClaw, Nanobot, GitHub Copilot, GitHub Models, GitClaw, and a plain system-prompt
target. We describe the repository-as-agent abstraction, the adapter model underlying
this portability, fourteen architectural patterns that emerge when agents are defined
git-natively (human-in-the-loop learning, agent versioning, branch-based deployment,
CI/CD for agents, segregation of duties, live memory, SkillsFlow, LLM Wiki, and others),
and a compliance model that maps directly to FINRA~3110/4511/2210,
SEC~Reg\,BI/Reg\,S-P/17a-4, Federal Reserve~SR\,11-7, and CFPB~Circular~2022-03.
Enforcement is achieved through \texttt{pre\_tool\_use} hooks with \texttt{fail\_open:
false} semantics, making policy a first-class, testable artifact rather than a runtime
wrapper. Spec v0.4 adds three concrete extensions to this surface: portable
\texttt{mcp\_servers} declarations that export to each runtime's native MCP
configuration, a \texttt{financial\_governance} block with spending caps and approval
thresholds for payment-capable agents, and an accepted RFC for an optional
Ed25519 cryptographic-identity layer.

OpenGAP is released under the MIT license. The reference implementation has 2{,}700+
GitHub stars, 15 adapters, 11 runtime runners, and ships provenance-signed as
\texttt{@open-gitagent/opengap} (v0.4.0) on npm. This paper consolidates the specification, the
architectural patterns, the portability model, and the compliance framing into a single
citable artifact, and sets out an agenda for conformance testing, formal
segregation-of-duties verification, and event-schema standardization for regulated
agent-to-agent interaction.
\end{abstract}

% ===============================================================
\section{Introduction}
% ===============================================================

The modern AI-agent ecosystem has two layers that are frequently conflated.

The first is the \emph{interface} layer: how an agent talks to tools, how two agents
exchange messages, how a user interface streams tokens back from a model. This layer is
well-served by existing standards. MCP~\cite{mcp2024} standardizes tool calls; Google's
A2A~\cite{a2a2024} standardizes agent-to-agent messaging; the Linux-Foundation ACP effort
and specifications such as AG-UI handle user-facing streams and event transport.

The second is the \emph{definition} layer: \emph{what the agent actually is}. Its
identity, its hard constraints, its permitted duties, its skills, the documents it may
consult, the memory it accumulates between sessions, the hooks that enforce its
boundaries, the regulatory artifacts that justify its deployment. Today this layer is
fragmented across framework-specific Python classes, hand-rolled YAML files,
dashboard-only configuration in vendor products, system prompts copied between Slack
channels, and screenshots pasted into compliance binders.

This fragmentation has three consequences. \textbf{(i)~Non-portability}: an agent
defined in CrewAI cannot be run in Claude Code without a rewrite.
\textbf{(ii)~Non-auditability}: a system prompt in a dashboard does not appear in
\texttt{git blame} and is not covered by SEC~17a-4 retention.
\textbf{(iii)~Non-governance}: segregation of duties, human-in-the-loop supervision,
and spending controls are wrapped around the agent by a runtime rather than being part
of the agent's own definition, so they can be bypassed by choosing a different
runtime.

We argue that the correct remedy is to standardize the definition layer as a git
repository, and we present \textbf{OpenGAP}\,---\,the GitAgentProtocol\,---\,to do
so. The contributions of this paper are:

\begin{enumerate}
  \item The \textbf{repository-as-agent abstraction} (\S\ref{sec:protocol}): a directory
  layout and manifest schema under which every aspect of an agent\,---\,identity,
  rules, duties, skills, tools, knowledge, memory, hooks, sub-agents, and compliance
  artifacts\,---\,is a file, under version control, diff-reviewable, and branch-able.
  \item The \textbf{adapter model} (\S\ref{sec:portability}): a systematic way to
  translate the canonical definition into fifteen different execution environments with
  a measurable fidelity profile per target, validated by the reference implementation
  \texttt{opengap}.
  \item A catalogue of \textbf{fourteen architectural patterns}
  (\S\ref{sec:patterns}) that emerge when agents are defined git-natively, including
  human-in-the-loop via branch~+~PR, agent versioning via git tags, segregation of
  duties via branch protection, CI/CD validation, branch-based deployment, live memory,
  SkillsFlow, and the LLM Wiki pattern.
  \item A \textbf{compliance-by-construction} model (\S\ref{sec:compliance}) that maps
  FINRA, SEC, Federal Reserve, and CFPB controls onto the git substrate and specifies
  enforcement through \texttt{pre\_tool\_use} hooks, yielding an audit trail that is
  structurally identical to \texttt{git log}.
\end{enumerate}

The remainder of the paper is organized as follows.
\S\ref{sec:related} positions OpenGAP against MCP, A2A, ADL-style capability
schemas, raw YAML configuration, and framework-native definitions.
\S\ref{sec:protocol} specifies the protocol.
\S\ref{sec:portability} describes the adapter model and summarizes the fidelity
profile across the fifteen current targets.
\S\ref{sec:patterns} catalogues the lifecycle patterns.
\S\ref{sec:compliance} develops the enterprise and regulatory story.
\S\ref{sec:impl} describes the reference implementation.
\S\ref{sec:cases} presents three case studies.
\S\ref{sec:eval} reports adoption signals and an export-fidelity evaluation.
\S\ref{sec:discuss} discusses limitations, a threat model, and governance.
\S\ref{sec:future} sets out future work, and \S\ref{sec:conclusion} concludes.

% ===============================================================
\section{Related Work}\label{sec:related}
% ===============================================================

\paragraph{Interface protocols: MCP, A2A, ACP, AG-UI.}
Anthropic's Model Context Protocol~\cite{mcp2024} specifies how a model host invokes
external tools and resources; Google's Agent-to-Agent protocol~\cite{a2a2024} specifies
how two autonomous agents negotiate capabilities and exchange task requests; the Linux
Foundation's ACP and AG-UI specifications cover agent-user interaction streams. All four
are interface specifications: they assume that an agent already exists and standardize
one edge of it. None of them specify the agent itself\,---\,its identity, its
constraints, its duties, or its regulatory artifacts. OpenGAP is complementary: a GAP
agent may expose MCP-compatible tools under \texttt{tools/}, declare A2A compatibility
in \texttt{agent.yaml} under \texttt{a2a:}, and stream via AG-UI at runtime, while the
GAP spec governs the definition that sits behind those interfaces.

\paragraph{Agent definition languages (ADLs).}
Several efforts propose single-file YAML or JSON schemas for describing an agent's
\emph{capabilities} for discovery or marketplace purposes. These resemble OpenAPI
specifications for agents. OpenGAP differs in two ways. First, the unit of description
is a directory, not a file: \texttt{SOUL.md}, \texttt{RULES.md}, \texttt{DUTIES.md},
\texttt{skills/*/SKILL.md}, \texttt{knowledge/}, \texttt{memory/}, and \texttt{hooks/}
are all first-class and all carry behavioral information that a single-file schema
cannot capture without becoming unwieldy. Second, a GAP repository is directly
\emph{runnable} and \emph{exportable}; it is not merely a capability advertisement.
The two approaches compose: a GAP \texttt{agent.yaml} may embed an ADL-style capability
manifest in its \texttt{tools} or \texttt{capabilities} block.

\paragraph{Framework-native definitions.}
LangChain, LangGraph, CrewAI, AutoGen, and similar frameworks define the agent as a
Python or TypeScript class. This gives full programming-language expressivity, but
yields three problems that OpenGAP addresses: (i)~the agent is not portable\,---\,a
CrewAI agent cannot be trivially executed under OpenAI Agents SDK; (ii)~the agent's
identity is not separable from its runtime, so a compliance reviewer must read code
rather than a manifest; and (iii)~version history is line-level rather than
semantic\,---\,\texttt{git log} shows diffs but not ``v2.1.0 changed the supervisory
policy.'' OpenGAP sits above frameworks: the canonical definition is the repository,
and the framework-native implementation is a derived artifact produced by
\texttt{opengap export}.

\paragraph{Raw YAML/JSON configuration.}
Many teams reach for a hand-rolled \texttt{config.yaml}. This offers flexibility but no
community-standard schema, no tooling layer, no reusable-skill ecosystem, and no
regulatory mapping. The OpenGAP specification (\S\ref{sec:protocol}) is what
\texttt{config.yaml} would be if teams agreed on one.

\paragraph{Summary.}
OpenGAP is the \emph{definition} standard to MCP's \emph{tool-call} standard and A2A's
\emph{agent-to-agent} standard. Table~\ref{tab:related} summarizes the layer each
addresses.

\begin{table}[h]
\centering
\footnotesize
\begin{tabular}{@{}lllll@{}}
\toprule
Protocol & Layer & Unit & Portable? & Compliance story \\
\midrule
MCP~\cite{mcp2024}         & Tool call      & JSON-RPC endpoint    & n/a          & out of scope \\
A2A~\cite{a2a2024}         & Agent-agent    & Message format       & n/a          & out of scope \\
ACP / AG-UI                & Agent-user     & Event stream         & n/a          & out of scope \\
ADL-style                  & Capability     & YAML/JSON file       & by adoption  & out of scope \\
Framework-native           & Implementation & Python/TS class      & no           & ad-hoc \\
Raw YAML                   & Config         & File                 & no           & none \\
\textbf{OpenGAP}          & \textbf{Definition} & \textbf{Git repo} & \textbf{yes (15 adapters)} & \textbf{first-class} \\
\bottomrule
\end{tabular}
\caption{OpenGAP occupies the definition layer; interface protocols are complementary.}
\label{tab:related}
\end{table}

% ===============================================================
\section{The OpenGAP Protocol}\label{sec:protocol}
% ===============================================================

\subsection{Design principles}

OpenGAP is grounded in four principles.

\begin{description}
  \item[P1. Git-native.] Version control, branching, diffing, tagging, and PR review are
  free operations. The agent definition is designed so that standard git workflows map
  cleanly onto the operations of the agent lifecycle (learning, deployment, rollback,
  review).
  \item[P2. Framework-agnostic.] The canonical definition does not presuppose a runtime.
  A deterministic export layer produces framework-specific artifacts on demand.
  \item[P3. Declarative over programmatic.] Wherever practicable, behavior is described
  in structured files (YAML, Markdown, JSON schemas) rather than code. Programmatic
  escape hatches exist (hook scripts, tool implementations), but the default is
  data.
  \item[P4. Composable.] Agents can extend, depend on, and delegate to other agents,
  each itself a git repository, pinned by version.
\end{description}

\subsection{The repository-as-agent abstraction}\label{sec:repo-as-agent}

A GAP agent is a directory. Figure~\ref{fig:layout} shows the canonical layout. Only
two files are strictly required: \texttt{agent.yaml}, which is validated against a
JSON~Schema~\cite{schema2020}, and \texttt{SOUL.md}, which is free-form Markdown.
All other files and directories are optional and additive.

\begin{figure}[t]
\begin{lstlisting}[basicstyle=\ttfamily\scriptsize]
my-agent/
|-- agent.yaml              # [REQUIRED] Manifest (validated)
|-- SOUL.md                 # [REQUIRED] Identity, personality, style
|-- RULES.md                # Hard constraints ("must never", "must always")
|-- DUTIES.md               # Segregation of duties policy / role boundaries
|-- AGENTS.md               # Framework-agnostic fallback instructions
|-- README.md               # Human documentation
|-- skills/                 # Reusable capability modules
|   `-- code-review/{SKILL.md, scripts/, references/, assets/, examples/}
|-- tools/                  # MCP-compatible tool schemas + impls
|-- knowledge/              # Reference documents (index.yaml + *.md|csv|pdf)
|-- memory/                 # Persistent cross-session memory (MEMORY.md, archive/)
|-- skillflows/             # Deterministic multi-step YAML procedures
|-- hooks/                  # Lifecycle handlers (hooks.yaml + scripts/)
|-- examples/               # Calibration interactions (few-shot)
|-- agents/                 # Sub-agent definitions (recursive)
|-- compliance/             # Regulatory artifacts
|   `-- {regulatory-map.yaml, audit-log.schema.json, risk-assessment.md, ...}
|-- config/                 # Environment-specific overrides (default.yaml, prod.yaml)
`-- .gitagent/              # Runtime state (gitignored)
\end{lstlisting}
\caption{Canonical GAP directory layout. Any directory not used is simply absent.}
\label{fig:layout}
\end{figure}

\subsection{Identity, rules, and duties}

Three files carry the agent's normative content.

\begin{itemize}
  \item \textbf{\texttt{SOUL.md}} is the agent's identity: its personality, its voice,
  the values it cares about, its communication style. It is short (typically under a
  page) and stable across releases. It is the first document a reviewer reads.
  \item \textbf{\texttt{RULES.md}} is the set of hard constraints: what the agent must
  always do, and what it must never do. RULES are tested by the \texttt{audit}
  command and referenced by enforcement hooks (\S\ref{sec:compliance}).
  \item \textbf{\texttt{DUTIES.md}} declares the agent's role within a multi-agent
  topology: what it is authorized to do, what it is explicitly forbidden from doing,
  and which roles it cannot hold simultaneously. This is the file that operationalizes
  segregation of duties for regulated deployments.
\end{itemize}

\subsection{Capabilities: skills, tools, MCP servers, knowledge, memory, hooks}

A GAP agent's capabilities are factored into six concerns.

\textbf{Skills} (\texttt{skills/<name>/SKILL.md}) are reusable capability modules:
each is a Markdown file with YAML frontmatter describing inputs, outputs, and
invocation conditions, plus optional scripts, references, assets, and examples.
Skills are installable from a shared registry; a single skill can be used by
multiple agents.

\textbf{Tools} (\texttt{tools/<name>.yaml}) are MCP-compatible schemas describing
an external function the agent may invoke, annotated with cost
(\texttt{none|low|medium|high}) and side-effect class. Implementations may live in
\texttt{tools/<name>.\{py,sh,js\}} or at an external endpoint.

\textbf{MCP servers} (\texttt{mcp\_servers} in \texttt{agent.yaml}) declare external
Model Context Protocol servers once and export them to each runtime's native
configuration. Each entry is either stdio-based (\texttt{command} + \texttt{args} +
\texttt{env}) or HTTP-based (\texttt{url} + \texttt{headers}), with \texttt{\$\{VAR\}}
interpolation for secrets. A single declaration renders to \texttt{.mcp.json} for
Claude Code, to the \texttt{mcpServers} block for Codex, to \texttt{.cursor/mcp.json}
for Cursor, and to a human-readable section for markdown-shaped targets. This extends
the ``define once, run anywhere'' property from the agent's own behavior to its
infrastructure integrations.

\textbf{Knowledge} (\texttt{knowledge/}) is reference material the agent may consult:
Markdown, CSV, PDF, or any readable format, plus an \texttt{index.yaml} that hints
retrieval keys. Knowledge is read-only from the agent's perspective.

\textbf{Memory} (\texttt{memory/}) is persistent cross-session state: a primary
\texttt{MEMORY.md} of at most 200 lines is always loaded; \texttt{memory/runtime/}
carries live state (daily log, current context); \texttt{memory/archive/} holds
historical snapshots that age out of the primary window. Memory is append-only by
convention and audit-friendly.

\textbf{Hooks} (\texttt{hooks/hooks.yaml}) register handlers for seven lifecycle
events: \texttt{on\_session\_start}, \texttt{pre\_tool\_use}, \texttt{post\_tool\_use},
\texttt{on\_error}, \texttt{user\_prompt\_submit}, \texttt{agent\_spawn}, and
\texttt{on\_session\_end}. Each hook has a \texttt{fail\_open} flag; setting
\texttt{fail\_open: false} on \texttt{pre\_tool\_use} upgrades a hook from an advisory
to an enforcement boundary (see \S\ref{sec:hook-enforcement}).

\subsection{Composition: extends, dependencies, sub-agents}

An agent may extend a parent (\texttt{extends: <git-url>}), shadowing only the files
it overrides. An agent may depend on peers (\texttt{dependencies:}) that are
shallow-cloned into specified mount paths at a pinned version. An agent may define
sub-agents in \texttt{agents/<name>/}, each itself a full GAP directory, and delegate
to them via a structured handoff that carries enforced duty scope.

\subsection{Formal manifest}

Appendix~\ref{app:schema} gives an excerpt of \texttt{agent-yaml.schema.json}, the
JSON Schema that \texttt{opengap validate} enforces on every \texttt{agent.yaml}. The
schema currently has ten files covering the agent manifest, hooks, hook I/O, tools,
skills, knowledge, memory, skillflows, configuration, and the marketplace descriptor.

% ===============================================================
\section{Portability: The Adapter Model}\label{sec:portability}
% ===============================================================

\subsection{Export-time versus runtime execution}

A GAP agent can be consumed in two ways. In \emph{export} mode, \texttt{opengap export
--format <target>} deterministically renders the canonical definition into a directory
or file in the target framework's native format (for example, \texttt{.cursor/rules/}
files for Cursor, \texttt{CLAUDE.md}~+~\texttt{.claude/} for Claude Code,
\texttt{AGENTS.md}~+~\texttt{opencode.json} for OpenCode). In \emph{runtime} mode,
\texttt{opengap run --adapter <target>} performs the equivalent export into a temporary
workspace and launches the target runtime over it, one-shot or interactively.

This separation matters. Export produces a reviewable artifact; runtime produces a
running process. A team that wants to own the execution environment picks export;
a team that wants convenience picks runtime.

\subsection{Adapter inventory}

Table~\ref{tab:adapters} enumerates the fifteen targets currently shipped. Each row
lists the adapter name, whether export and/or runtime are supported, the external
dependency required, and the target's native input format.

\begin{table}[t]
\centering
\footnotesize
\begin{tabular}{@{}lccll@{}}
\toprule
Adapter & Export & Runtime & Required & Target native format \\
\midrule
\texttt{claude-code}    & \checkmark & \checkmark & Claude Code CLI               & \texttt{CLAUDE.md} + \texttt{.claude/} \\
\texttt{cursor}         & \checkmark & -          & Cursor IDE                    & \texttt{.cursor/rules/*.mdc} \\
\texttt{codex}          & \checkmark & -          & OpenAI Codex CLI              & \texttt{AGENTS.md} + config \\
\texttt{gemini}         & \checkmark & \checkmark & Google Gemini CLI             & \texttt{GEMINI.md} + \texttt{settings.json} \\
\texttt{opencode}       & \checkmark & \checkmark & OpenCode CLI                  & \texttt{AGENTS.md} + \texttt{opencode.json} \\
\texttt{kiro}           & \checkmark & -          & Kiro CLI                      & \texttt{.kiro/agents/*.json} \\
\texttt{openai}         & \checkmark & \checkmark & OpenAI Agents SDK             & Python script \\
\texttt{crewai}         & \checkmark & \checkmark & CrewAI CLI                    & YAML config \\
\texttt{lyzr}           & \checkmark & \checkmark & Lyzr Studio                   & Lyzr agent definition \\
\texttt{openclaw}       & \checkmark & \checkmark & OpenClaw CLI                  & Workspace layout \\
\texttt{nanobot}        & \checkmark & \checkmark & Nanobot CLI                   & Nanobot config \\
\texttt{copilot}        & \checkmark & -          & GitHub Copilot                & \texttt{.github/copilot-*} \\
\texttt{github}         & \checkmark & \checkmark & GitHub Models                 & GitHub Actions workflow \\
\texttt{gitclaw}        & \checkmark & \checkmark & GitClaw                       & GitClaw workspace \\
\texttt{system-prompt}  & \checkmark & -          & any LLM                       & Concatenated text \\
\bottomrule
\end{tabular}
\caption{The fifteen adapters shipped in the reference implementation. All export
adapters are deterministic: the same GAP directory yields the same output bytes.}
\label{tab:adapters}
\end{table}

\subsection{Fidelity profile}

Not every adapter preserves every element of a GAP definition. Some targets have no
notion of sub-agents; some cannot carry hooks; some accept only a single system-prompt
blob. We define the \emph{fidelity profile} of an adapter as the set of GAP
elements it faithfully represents in its target format. Table~\ref{tab:fidelity}
summarizes the profile at a coarse grain; the full matrix (Appendix~\ref{app:fidelity})
is regenerated by running \texttt{opengap export} against the reference agents
\texttt{examples/standard/} and \texttt{examples/full/} and mechanically diffing the
output for the presence of each element.

\begin{table}[h]
\centering
\footnotesize
\begin{tabular}{@{}lccccccccc@{}}
\toprule
Adapter & \textsc{soul} & \textsc{rules} & \textsc{duties} & skills & tools & hooks & memory & sub-ag. & compl. \\
\midrule
\texttt{claude-code}   & F & F & F & F & F & F & F & F & F \\
\texttt{cursor}        & F & F & P & F & F & - & - & - & P \\
\texttt{opencode}      & F & F & P & F & F & F & P & - & F \\
\texttt{gemini}        & F & F & P & F & F & - & P & - & P \\
\texttt{codex}         & F & F & P & P & F & - & - & - & P \\
\texttt{kiro}          & F & F & P & F & F & F & - & - & P \\
\texttt{openai}        & F & F & P & F & F & - & - & - & - \\
\texttt{crewai}        & F & F & P & F & F & - & - & F & - \\
\texttt{lyzr}          & F & F & P & F & F & - & P & - & P \\
\texttt{openclaw}      & F & F & P & F & F & P & F & F & F \\
\texttt{nanobot}       & F & F & P & F & F & P & P & - & - \\
\texttt{system-prompt} & F & F & P & P & - & - & - & - & - \\
\bottomrule
\end{tabular}
\caption{Abbreviated fidelity profile. \textbf{F} = fully preserved, \textbf{P} =
partially (text only, semantics dropped), \textbf{-} = not representable. Rows for
\texttt{copilot}, \texttt{github}, \texttt{gitclaw} omitted for space; see
Appendix~\ref{app:fidelity}.}
\label{tab:fidelity}
\end{table}

The matrix clarifies a point often glossed over in portability claims: complete
bidirectional interop is impossible because target frameworks differ in expressive
power. What is possible is \emph{deterministic, lossy-where-necessary, documented}
export\,---\,which is what OpenGAP provides.

% ===============================================================
\section{Lifecycle Patterns}\label{sec:patterns}
% ===============================================================

Defining agents git-natively yields architectural patterns that have no clean analogue
in framework-native or dashboard-configured agents. We catalogue fourteen such patterns
here; each has a corresponding diagram in \texttt{patterns/} in the reference
repository.

\paragraph{P1: Human-in-the-Loop via Branch~+~PR.}
When an agent proposes a new skill, a new rule, or a change to its own memory, it
opens a branch and a pull request. A human reviewer approves the PR before it merges
into \texttt{main}. Reinforcement-learning-style updates are thereby subject to
code-review semantics: diffs, comments, blocking approvals, revert via
\texttt{git~revert}.

\paragraph{P2: Agent Versioning via Git Tags.}
Semantic version tags on the agent repository act as the deployment unit. A runtime
consumes \texttt{github.com/org/agent@v2.1.0}. Rollback is \texttt{git checkout v2.0.3}
rather than a dashboard operation. Undo history is the commit graph.

\paragraph{P3: Segregation of Duties via Branch Protection.}
\texttt{DUTIES.md} declares roles; the agent cannot approve its own work because it
cannot merge its own branch. This is enforced by branch-protection rules on the hosting
platform, not by the agent's internal logic\,---\,a structural guarantee rather than a
policy assertion.

\paragraph{P4: CI/CD for Agents.}
\texttt{opengap validate --compliance} runs in continuous integration on every push.
A failing agent manifest blocks the merge. Agent quality is treated as code quality.

\paragraph{P5: Branch-Based Deployment.}
\texttt{dev}, \texttt{staging}, and \texttt{main} branches correspond to deployment
environments. An agent is promoted by merging. Environment-specific configuration
lives in \texttt{config/<env>.yaml} and is selected at runtime.

\paragraph{P6: Live Agent Memory.}
\texttt{memory/MEMORY.md} is the 200-line cross-session memory; \texttt{memory/runtime/}
holds a daily log and the current task context. Memory writes are commits, preserving
the history of what the agent knew when.

\paragraph{P7: Shared Context \& Skills.}
When multiple agents share a working directory, root-level \texttt{skills/},
\texttt{tools/}, and \texttt{knowledge/} are common across sub-agents defined in
\texttt{agents/<name>/}. This yields a natural monorepo pattern without repository
sprawl.

\paragraph{P8: Agent Forking \& Remixing.}
\texttt{extends: <git-url>} lets an agent fork a parent and shadow only the files it
wishes to change. Open-source agents become a remix culture; vendor-specific overrides
live in forks.

\paragraph{P9: Knowledge Tree.}
\texttt{knowledge/index.yaml} describes a hierarchical set of documents with
retrieval hints; embeddings and search indices are derived artifacts
materialized in \texttt{.gitagent/cache/}.

\paragraph{P10: Agent Diff \& Audit Trail.}
\texttt{git diff} between two tags is a diff of the agent's identity, rules, duties,
and capabilities. Reviewers see exactly what changed between \texttt{v2.1.0} and
\texttt{v2.2.0}\,---\,a property worth a great deal in regulated deployments.

\paragraph{P11: Tagged Releases.}
Tags are signed and optionally immutable. \texttt{v2025-Q1-audit} is the signed-off
state of an agent at the close of a quarter. Pinning consumers to a tag is standard
practice.

\paragraph{P12: Secret Management via \texttt{.env}.}
Secrets are never committed. \texttt{.env} is in \texttt{.gitignore}; the spec
declares required environment variables in \texttt{agent.yaml} under
\texttt{env\_vars:}, and the reference implementation warns on missing variables
before runtime.

\paragraph{P13: Lifecycle Hooks.}
\texttt{bootstrap.md} runs at session start; \texttt{teardown.md} at session end; the
seven-event hook system covers pre/post tool use, errors, prompt submission, and
sub-agent spawn. Hooks are first-class, versioned alongside the agent.

\paragraph{P14: SkillsFlow.}
Deterministic multi-step workflows are described in \texttt{skillflows/*.yaml} with
explicit \texttt{depends\_on} ordering and \texttt{\${{...}}} variable substitution.
Where a skill is probabilistic (LLM reasoning), a skillflow is deterministic (a DAG
over skills). SkillsFlows are the GAP counterpart to LangGraph's \texttt{StateGraph}.

\vspace{0.5em}
These patterns are not mere conveniences. Several\,---\,P1, P3, P4, P10\,---\,are
structural preconditions for deploying agents under FINRA 3110 supervision or
SR~11-7 model-risk management.

% ===============================================================
\section{Enterprise and Regulatory}\label{sec:compliance}
% ===============================================================

\subsection{Compliance as a first-class spec element}

Most agent frameworks treat compliance as a runtime wrapper: a governance service
intercepts tool calls, a policy engine audits responses, an SRE team files retention
policies manually. The agent itself is unaware of its regulatory context. This is
brittle\,---\,swapping the runtime loses the controls\,---\,and opaque\,---\,a
compliance reviewer cannot read the agent's definition and determine what regime
applies.

OpenGAP inverts this. A \texttt{compliance:} block in \texttt{agent.yaml} and an
optional \texttt{compliance/} directory carry the regulatory artifacts directly:

\begin{lstlisting}[language=yamlex]
compliance:
  framework: FINRA
  risk_tier: high              # SR 11-7 model-risk tier
  rules:
    - rule: "3110"             # FINRA supervision
      description: All outputs reviewed before client delivery
    - rule: "4511"             # FINRA books-and-records
      description: Immutable audit trail of all interactions
  communications:
    fair_balanced: true        # FINRA 2210
    no_misleading: true
  data_governance:
    pii_handling: redact       # Reg S-P
    retention_years: 7         # SEC 17a-4
  supervision:
    human_in_the_loop: always
  segregation_of_duties:
    roles:
      - id: analyst
        permissions: [draft, revise]
      - id: reviewer
        permissions: [review, approve, reject]
    conflicts:
      - [analyst, reviewer]
    enforcement: strict
\end{lstlisting}

\subsection{Regulatory mapping}

Table~\ref{tab:regs} maps GAP fields onto the specific regulations they support. This
is not a legal opinion; it is a structural mapping that auditors use to demonstrate
control coverage.

\begin{table}[h]
\centering
\footnotesize
\begin{tabular}{@{}lll@{}}
\toprule
Regulation & Control & GAP field(s) \\
\midrule
FINRA~3110          & Supervision           & \texttt{supervision.human\_in\_the\_loop}, hook PR review \\
FINRA~3120          & Compliance testing    & \texttt{compliance/validation-schedule.yaml} \\
FINRA~4511          & Books and records     & \texttt{compliance.recordkeeping.retention\_years} \\
FINRA~2210          & Communications        & \texttt{communications.fair\_balanced} \\
Reg~Notice~24-09    & GenAI/LLM application & entire \texttt{compliance:} block \\
SEC~Reg~BI          & Best interest         & \texttt{RULES.md} best-interest constraints \\
SEC~Reg~S-P         & Privacy / PII         & \texttt{data\_governance.pii\_handling} \\
SEC~17a-4           & Electronic records    & \texttt{retention\_years: 7}, git history \\
Federal~Res.~SR~11-7& Model risk management & \texttt{risk\_tier}, \texttt{compliance/risk-assessment.md} \\
Federal~Res.~SR~23-4& AI risk management    & \texttt{compliance/validation-schedule.yaml} \\
CFPB~Circular~2022-03 & Adverse-action notices & \texttt{RULES.md} explainability requirements \\
\bottomrule
\end{tabular}
\caption{GAP fields mapped onto US financial regulations and guidance.}
\label{tab:regs}
\end{table}

\subsection{Hook-based enforcement}\label{sec:hook-enforcement}

Declaring a policy is not enforcing it. Enforcement happens at \texttt{pre\_tool\_use}:
before any tool call, the registered pre-tool hook fires, receives the tool name and
arguments on \texttt{stdin}, and exits \texttt{0} to allow or non-zero to block. A
hook configured with \texttt{fail\_open: false} blocks the tool call on error, timeout,
or non-zero exit. This is the mechanism by which spending caps, PII redaction, and
segregation-of-duties conflicts are structurally enforced.

\subsection{Financial governance}\label{sec:fingov}

As of spec v0.4, the \texttt{compliance.financial\_governance} block is a shipped,
schema-validated feature for payment-capable agents. It declares spending limits in
absolute minor units (cents), an approval threshold, allowed and blocked categories,
and a named \emph{firewall} identifier\,---\,never an endpoint URL, keeping the spec
vendor-neutral:

\begin{lstlisting}[language=yamlex]
compliance:
  risk_tier: high
  financial_governance:
    enabled: true
    firewall: valkurai          # named identifier, not an endpoint
    spending:
      max_per_transaction_cents: 5000     # $50.00 hard cap
      max_monthly_cents: 100000           # $1,000.00 cumulative
      currency: AUD                       # ISO 4217
      allowed_categories: [software, compute, api_services]
      blocked_categories: [gambling, crypto, unknown]
    approval:
      require_above_cents: 2000           # human approval above $20.00
      timeout_minutes: 60
      auto_deny_on_timeout: true
\end{lstlisting}

The block is declarative; enforcement is a \texttt{pre\_tool\_use} hook that reads
\texttt{financial\_governance.spending.max\_per\_transaction\_cents}, parses the pending
payment from tool input, and blocks the call if the cap is exceeded. \texttt{opengap
validate} warns when a \texttt{risk\_tier: high} agent declares financial tools but no
\texttt{financial\_governance} block, and rejects a \texttt{firewall} value that looks
like a URL. The hook lives in the agent repository under version control, reviewable in
the same pull request as the policy it enforces\,---\,a locality that runtime-only
policy engines cannot provide. The event schema for cross-runtime payment interop is
deliberately deferred (\S\ref{sec:future}) until two or more enforcement implementations
exist; standardizing it earlier would standardize imaginary interop.

\subsection{Audit trail is \texttt{git log}}

Every commit to the agent repository is an audit-eligible event. \texttt{git log}
gives a timestamped, attributed, immutable sequence. \texttt{git blame} answers
``who changed this rule, when, and why'' in constant time. No separate audit store is
needed. For regulated domains, an organization pins its agents to signed tags, writes
post-hoc review findings as PR comments, and archives the repository on
retention-compliant storage. The compliance property is structural rather than
bolted-on. See \cite{gitclose} for a worked example of this pattern applied to the
monthly financial close.

\subsection{Threat model: what OpenGAP does not prevent}\label{sec:threat}

Honesty about limitations is as valuable as a feature list. OpenGAP does \emph{not}
prevent the following:

\begin{enumerate}
  \item \textbf{A runtime that ignores hooks.} If a target framework does not implement
  \texttt{pre\_tool\_use} hooks, enforcement silently degrades to advisory. The export
  fidelity matrix (Table~\ref{tab:fidelity}) makes this explicit per target.
  \item \textbf{An adversary with commit rights.} Branch protection defeats accidents;
  it does not defeat insiders. Signed commits and code-owner reviews mitigate but do
  not eliminate this.
  \item \textbf{Prompt injection through knowledge.} Documents in \texttt{knowledge/}
  may carry payloads that override rules. Defensive measures (content filters,
  provenance checks) are out of scope for the spec.
  \item \textbf{Model capability drift.} A regulatory statement such as
  ``human\_in\_the\_loop: always'' does not verify that the human actually reviewed the
  output; it only requires that the runtime offer a review gate.
  \item \textbf{Extrinsic compliance obligations.} GAP maps structural controls onto
  git workflows. Jurisdiction-specific duties (MiCA, AI Act, HIPAA specifics) require
  local legal review; GAP provides the substrate, not the opinion.
\end{enumerate}

\subsection{Cryptographic identity (RFC, optional)}\label{sec:identity}

The threat-model item on insiders motivates an optional cryptographic identity layer,
accepted as an RFC (\texttt{spec/rfcs/identity.md}) in spec v0.4 and slated as an
optional \texttt{identity} block on \texttt{agent.yaml}. It distinguishes two concerns
the literature often conflates: \emph{provenance} (this manifest at this commit was
authored by the holder of key $X$\,---\,already solvable with signed tags or sigstore,
no spec change needed) and \emph{runtime delegation} (the running agent producing this
output acts on behalf of parent agent $Y$, with scope $Z$, signed by $Y$'s key, not yet
revoked\,---\,the genuine gap). The RFC binds a manifest to an Ed25519 public key, with
an optional \texttt{passport\_uri} pointing at a richer identity document for scoped
delegation and revocation, and reserves \texttt{signatures.<scope>} for manifest
signatures. Verification semantics live in the spec; enforcement (refuse to load,
sandbox, log-and-continue) is left to the runtime. The layer is fully optional;
manifests without it are unchanged. This closes the loop on inter-organization
delegation and the regulated-runtime check ``prove the agent that produced this output
was the agent in the manifest, with authority that had not been revoked.''

% ===============================================================
\section{Reference Implementation}\label{sec:impl}
% ===============================================================

The reference implementation is \texttt{opengap}, a TypeScript CLI published to npm as
\texttt{@open-gitagent/opengap} (scoped, provenance-signed via GitHub Actions OIDC) at
v0.4.0. The \texttt{gitagent} command is installed as a backward-compatibility alias
for the same binary. The repository is \texttt{github.com/open-gitagent/opengap}; the
license is MIT. (The project was originally named \texttt{gitagent}, briefly
\texttt{gapman}; the unscoped root name \texttt{opengap} is unavailable on npm under its
package-similarity policy, so the scoped name is canonical.)

\texttt{opengap} exposes the following verbs: \texttt{init} (scaffold from templates
\texttt{minimal}, \texttt{standard}, \texttt{full}, \texttt{llm-wiki}),
\texttt{validate} (JSON-schema validation and optional compliance audit),
\texttt{info} (summarize an agent), \texttt{export} (render to a target framework),
\texttt{import} (convert from Claude, Cursor, CrewAI, OpenCode, Gemini, or Codex
format), \texttt{install} (resolve git-based dependencies), \texttt{audit}
(compliance report), \texttt{skills} (search, install, list, inspect),
\texttt{run} (one-shot or interactive execution under an adapter),
\texttt{registry} (browse the community registry), and \texttt{lyzr} (Lyzr-Studio
integration). The CLI's help surface and semantics track the specification strictly.

Adapters and runners are separated: \texttt{src/adapters/} contains deterministic
renderers, while \texttt{src/runners/} contains the code that provisions a workspace
and launches the target runtime. Adapters are pure functions of the agent directory;
runners additionally manage processes, credentials, and interactive streams.

% ===============================================================
\section{Case Studies}\label{sec:cases}
% ===============================================================

\subsection{NVIDIA Deep Researcher}

\texttt{examples/nvidia-deep-researcher/} is a three-agent hierarchy modeled on NVIDIA
AIQ's deep-research workflow\,---\,an orchestrator, a planner, and a researcher
sub-agent with three skills (paper search, web search, knowledge retrieval). The GAP
version captures the SOD hierarchy in \texttt{DUTIES.md}
(\texttt{segregation\_of\_duties.conflicts: [[orchestrator, researcher]]}), making
the structural guarantee that no single agent may both plan and execute the research
step. Exported to Claude Code, the hierarchy becomes a parent agent with two
sub-agents; exported to CrewAI, the same hierarchy becomes three crew members with a
task graph. The identity, duties, and skills are unchanged; only the native encoding
differs.

\subsection{Full-compliance financial agent}

\texttt{examples/full/} is a production-grade compliance agent. It includes a
\texttt{compliance/regulatory-map.yaml} mapping 12 specific rules to 12 controls,
a \texttt{compliance/risk-assessment.md} justifying a \texttt{risk\_tier: high}
classification under SR~11-7, and a \texttt{compliance/validation-schedule.yaml}
specifying annual bias testing and quarterly parallel testing. A
\texttt{hooks/pre\_tool\_use.sh} with \texttt{fail\_open: false} blocks tool
calls that violate the spending cap. The agent is not a toy: every file is
what a regulated team would produce, and the total diff between v1.0.0 and
v1.1.0 of the agent is a coherent, reviewable PR.

\subsection{LLM Wiki}

\texttt{examples/llm-wiki/} implements the knowledge-compile pattern described by
Karpathy~\cite{karpathy_llmwiki}: \texttt{knowledge/} holds raw source material,
\texttt{memory/wiki/} holds compiled wiki pages, and three
skills\,---\,\texttt{wiki-ingest}, \texttt{wiki-query}, \texttt{wiki-lint}\,---\,cover
the life cycle of the wiki. A \texttt{memory/log.md} operation log records every
ingestion and query, yielding a complete provenance trail. The case study
demonstrates that a substantive cognitive pattern from outside the GAP community
fits inside the protocol without spec changes.

% ===============================================================
\section{Evaluation}\label{sec:eval}
% ===============================================================

We report evaluation along three axes: adoption, fidelity, and qualitative comparison.

\subsection{Adoption signals}

As of \today, the reference repository \texttt{open-gitagent/opengap} has
\textbf{2{,}700+} GitHub stars and \textbf{15 export adapters} plus \textbf{11 runtime
runners}. Five adapters (Cursor, Kiro, Codex, Gemini, OpenCode) were contributed via
pull request by external authors, and three spec features now in \texttt{main}\,---\,
portable \texttt{mcp\_servers} (\S\ref{sec:protocol}), \texttt{financial\_governance}
(\S\ref{sec:fingov}), and the cryptographic-identity RFC (\S\ref{sec:identity})\,---\,
originated as community RFCs and pull requests. The package ships provenance-signed as
\texttt{@open-gitagent/opengap} v0.4.0; continuous integration builds on Node 18, 20,
and 22 and validates the bundled example agents on every push. Adapter and spec
contributions from external developers are the strongest signal that the protocol is
perceived as a neutral substrate rather than a vendor product\,---\,no contributor has
standing to force a breaking change in favor of their framework. These numbers are to be
regenerated from the GitHub API at the time of arXiv submission; we report them here as
a floor.

\subsection{Export fidelity}

The fidelity profile (Table~\ref{tab:fidelity}, full matrix in
Appendix~\ref{app:fidelity}) is derived by running
\texttt{opengap export --format <target>} over the reference agents
\texttt{examples/standard/} and \texttt{examples/full/} and diffing the output for the
presence of each GAP element. Two headline observations:

\begin{enumerate}
  \item \texttt{claude-code} is the only adapter that achieves full fidelity across
  all nine tested elements. This reflects that GAP and Claude Code were co-designed:
  GAP's directory layout is close to Claude Code's native \texttt{CLAUDE.md} +
  \texttt{.claude/} convention.
  \item Hooks and sub-agents are the most commonly lost elements on export.
  \texttt{cursor}, \texttt{gemini}, \texttt{codex}, \texttt{openai}, and
  \texttt{system-prompt} all drop hooks entirely. Policy enforcement via hooks is
  therefore a \emph{GAP-native} guarantee and a \emph{per-target} degradation; teams
  requiring enforcement should verify that their target runtime executes the hook
  layer.
\end{enumerate}

\subsection{Qualitative comparison}

Table~\ref{tab:qualitative} extends Table~\ref{tab:related} with feature-level
comparison. The contrast is strongest in compliance and portability axes and is
indistinguishable from the status-quo on fundamental interface concerns (tool calls,
streaming, messaging)\,---\,which is as intended: OpenGAP complements MCP/A2A rather
than competing with them.

\begin{table}[h]
\centering
\footnotesize
\begin{tabular}{@{}lccccc@{}}
\toprule
Capability & GAP & MCP & A2A & ADL-style & Framework-native \\
\midrule
Defines the agent's identity      & \checkmark & -          & -          & partial    & \checkmark \\
Versioned via git tags            & \checkmark & -          & -          & -          & partial \\
Portable across frameworks        & \checkmark (15) & -     & -          & by adoption & - \\
Compliance fields first-class     & \checkmark & -          & -          & -          & - \\
SOD structural guarantee          & \checkmark & -          & -          & -          & ad-hoc \\
Enforcement hook layer            & \checkmark & -          & -          & -          & ad-hoc \\
Runnable from repo URL            & \checkmark & -          & -          & -          & - \\
Tool-call format                  & uses MCP   & \checkmark & -          & -          & framework \\
Agent-to-agent messaging          & uses A2A   & -          & \checkmark & -          & framework \\
\bottomrule
\end{tabular}
\caption{Feature-level comparison. \checkmark = first-class, partial = supported but
ad-hoc, - = out of scope or absent.}
\label{tab:qualitative}
\end{table}

% ===============================================================
\section{Discussion}\label{sec:discuss}
% ===============================================================

\paragraph{Limitations.}
The specification is at v0.1.0. Several elements are deliberately under-formalized:
memory archival policy, skill-registry trust model, and the interaction between
\texttt{config/<env>.yaml} and \texttt{hooks/} in multi-tenant runtimes. A conformance
test suite does not yet exist; today's adapters are validated by unit tests and by
community use, not by a formal conformance profile. Formal semantics for
segregation-of-duties (\S\ref{sec:patterns} P3) are described in natural language
rather than verified by a model checker. These are all addressed in future work
(\S\ref{sec:future}).

\paragraph{Vendor-capture risk and governance.}
An open protocol with a dominant reference implementation is one bug report away from
becoming a de-facto vendor standard. OpenGAP addresses this through three structural
choices: (i)~the spec and the reference implementation are in the same repository,
under the same license (MIT), under the same MIT-licensed working-group
governance\,---\,no silent re-licensing; (ii)~adapter contributions from external
authors for Cursor, Kiro, Codex, Gemini, OpenCode, Copilot, and GitHub Models
demonstrate plural interest and reduce the chance that any one target dictates
changes; (iii)~named implementations (for example, proposed financial-governance
firewalls) are listed in docs as \emph{example} implementations, never as \emph{the
reference}, to avoid privileging a single vendor in spec text. The working group is
expanding contributor representation as adoption grows.

\paragraph{Why a protocol and not a framework.}
A framework solves problems for the framework's users. A protocol creates a substrate
that multiple competing products can build on. OpenGAP chooses the protocol
path\,---\,the reference implementation (\texttt{opengap}) is one of potentially many,
and the spec has deliberate latitude so that vendors can add value within target
adapters without needing to modify the canonical definition.

% ===============================================================
\section{Future Work}\label{sec:future}
% ===============================================================

\begin{enumerate}
  \item \textbf{Conformance test suite.} A portable, language-agnostic test suite that
  a prospective adapter implementer can run to certify fidelity claims. This is the
  single highest-leverage next artifact.
  \item \textbf{Formal SOD semantics.} Express \texttt{DUTIES.md} conflict rules as a
  set of logical assertions, and verify that an agent's tool-call trace satisfies
  them. A model checker such as TLA$^+$ or Alloy is a natural fit.
  \item \textbf{Financial-governance Phase 2/3.} The declarative
  \texttt{financial\_governance} block shipped in v0.4 (\S\ref{sec:fingov}); Phase 2 is
  a reference \texttt{pre\_tool\_use} enforcement hook, and Phase 3 is a
  \texttt{payment\_required}~/ \texttt{payment\_approval}~/ \texttt{payment\_receipt}
  event schema for cross-runtime interop. Phase 3 waits until two or more enforcement
  implementations exist; standardizing earlier risks codifying imaginary
  interop~\cite{rfc_38}.
  \item \textbf{Identity-layer implementation.} The cryptographic-identity RFC
  (\S\ref{sec:identity}) is accepted; the next step is the \texttt{identity} block in
  \texttt{agent-yaml.schema.json}, a reference verifier, and a delegation-chain example.
  \item \textbf{Deeper A2A integration.} GAP currently declares A2A compatibility in
  \texttt{agent.yaml} but does not ship an A2A server adapter. A reference A2A
  adapter would close the loop between agent-to-agent interfaces and GAP definitions.
  \item \textbf{Empirical user study.} A controlled study of team velocity when
  defining, porting, and auditing agents under GAP versus framework-native
  equivalents. This is the evaluation that a full systems-venue submission would
  require.
\end{enumerate}

% ===============================================================
\section{Conclusion}\label{sec:conclusion}
% ===============================================================

Agents are now a production software category. The definition layer\,---\,what an
agent is, what it is permitted to do, what it is accountable for, and how it is
versioned over time\,---\,has until now been served by framework-native code,
hand-rolled YAML, or dashboard configuration, none of which are portable, auditable,
or governance-friendly.

OpenGAP standardizes this layer as a git repository. A single canonical definition
exports to fifteen execution environments with a documented fidelity profile; fourteen
lifecycle patterns emerge naturally from the git substrate (human-in-the-loop,
versioning, segregation of duties, CI/CD, branch-based deployment, live memory,
SkillsFlow, LLM Wiki, and others); and compliance becomes a first-class spec element
mapped directly onto FINRA, SEC, Federal Reserve, and CFPB controls, with enforcement
implemented through hooks that are themselves under version control.

The protocol is MIT-licensed, the reference implementation is 2{,}700+-star
community-maintained, and the full specification, patterns, and adapters ship together
in one repository. We invite the research and practitioner community to contribute
adapters, conformance tests, case studies, and critical review.

% ===============================================================
\section*{Acknowledgments}
% ===============================================================
\noindent
The OpenGAP working group thanks contributors to the reference implementation and to
the Cursor, Kiro, Codex, Gemini, OpenCode, and Copilot adapters. Early design
conversations with the MCP and A2A communities sharpened the scoping of OpenGAP as a
complementary definition layer. All remaining mistakes are our own.

% ===============================================================
\bibliographystyle{plain}
\bibliography{bibliography}
% ===============================================================

\newpage
\appendix

% ===============================================================
\section{agent.yaml Schema (Excerpt)}\label{app:schema}
% ===============================================================

\begin{lstlisting}[basicstyle=\ttfamily\scriptsize]
{
  "$id": "https://gitagent.sh/spec/v0.1.0/agent-yaml.schema.json",
  "type": "object",
  "required": ["name", "version", "description"],
  "properties": {
    "spec_version":  { "type": "string", "pattern": "^\\d+\\.\\d+\\.\\d+$" },
    "name":          { "type": "string", "pattern": "^[a-z][a-z0-9-]*$" },
    "version":       { "type": "string", "pattern": "^\\d+\\.\\d+\\.\\d+(?:[-+].+)?$" },
    "description":   { "type": "string" },
    "model": {
      "type": "object",
      "properties": {
        "preferred": { "type": "string" },
        "fallback":  { "type": "array", "items": { "type": "string" } },
        "constraints": {
          "type": "object",
          "properties": {
            "temperature":   { "type": "number", "minimum": 0, "maximum": 2 },
            "max_tokens":    { "type": "integer", "minimum": 1 }
          }
        }
      }
    },
    "compliance": {
      "type": "object",
      "properties": {
        "framework":           { "type": "string" },
        "risk_tier":           { "enum": ["low", "medium", "high", "critical"] },
        "supervision":         { "type": "object" },
        "segregation_of_duties": { "type": "object" },
        "recordkeeping":       { "type": "object" }
      }
    },
    "dependencies": { "type": "array" },
    "extends":      { "type": "string" },
    "a2a":          { "type": "array", "items": { "type": "string" } }
  }
}
\end{lstlisting}

% ===============================================================
\section{Full Fidelity Matrix}\label{app:fidelity}
% ===============================================================

\begin{table}[h]
\centering
\scriptsize
\begin{tabular}{@{}lccccccccc@{}}
\toprule
Adapter & \textsc{soul} & \textsc{rules} & \textsc{duties} & skills & tools & hooks & memory & sub-ag. & compl. \\
\midrule
\texttt{claude-code}     & F & F & F & F & F & F & F & F & F \\
\texttt{cursor}          & F & F & P & F & F & - & - & - & P \\
\texttt{codex}           & F & F & P & P & F & - & - & - & P \\
\texttt{gemini}          & F & F & P & F & F & - & P & - & P \\
\texttt{opencode}        & F & F & P & F & F & F & P & - & F \\
\texttt{kiro}            & F & F & P & F & F & F & - & - & P \\
\texttt{openai}          & F & F & P & F & F & - & - & - & - \\
\texttt{crewai}          & F & F & P & F & F & - & - & F & - \\
\texttt{lyzr}            & F & F & P & F & F & - & P & - & P \\
\texttt{openclaw}        & F & F & P & F & F & P & F & F & F \\
\texttt{nanobot}         & F & F & P & F & F & P & P & - & - \\
\texttt{copilot}         & F & F & P & P & P & - & - & - & - \\
\texttt{github}          & F & F & P & P & F & - & - & - & P \\
\texttt{gitclaw}         & F & F & F & F & F & F & F & F & F \\
\texttt{system-prompt}   & F & F & P & P & - & - & - & - & - \\
\bottomrule
\end{tabular}
\caption{Full fidelity matrix for the fifteen shipped adapters. \textbf{F} = fully
preserved, \textbf{P} = partially preserved (text only or structure lost), \textbf{-}
= not representable in the target format.}
\end{table}

\end{document}
